\section{A conditional central limit argument and uniform Gamma}
\label{u:sec:gamma}

Throughout the construction of uniform property $\Gamma$, let $A$
satisfy the hypotheses of Theorem~\ref{u:thm:main} with
$T(A)\ne\varnothing$, and retain $D$ and $\Lambda$ from
Section~\ref{u:sec:preliminaries}.  We turn the bounded moment estimates
of Section~\ref{u:sec:averaging} into central projections that split
every coefficient trace equally.  A common atomless limiting
distribution will let us use the same scalar cutoff functions for
every trace.  We first transfer the local estimates to arbitrary
norm-separable coefficient algebras in $D$, then adjoin commuting
variables successively.  Their normalized sums will have Gaussian limiting
characteristic functions, and equal-measure Gaussian intervals will
give the required projections.

\subsection{Transfer to separable coefficient algebras}

The local conclusion of Proposition~\ref{u:prop:local-moments} applies
to coordinate lifts of elements of any norm-separable subalgebra of
$D$.

\begin{proposition}[Relative moment estimates]
\label{u:prop:relative-moments}
Let $C\subset D$ be a unital norm-separable $C^*$-subalgebra and fix
$R\geq1$.  There is a selfadjoint element $d\in D\cap C'$ with
$\norm{d}\leq R$ such that, for every $b\in C$ with $\norm{b}\leq1$ and
every $\tau\in\Lambda$,
\begin{equation}
\label{u:eq:relative-moments}
  |\tau(bd)|\leq\frac{3}{R^3},\qquad
  |\tau(b(d^2-1))|\leq\frac{3}{R^2},\qquad
  \tau(d^4)\leq3.
\end{equation}
\end{proposition}

\begin{proof}
Choose a norm-dense sequence $(c_k)_{k\geq1}$ in the unit ball of $C$
and a norm-dense sequence $(x_k)_{k\geq1}$ in its selfadjoint unit ball.
Choose lifts $(c_{k,l})_{l\geq1}$ consisting of contractions in $A$ and
lifts $(x_{k,l})_{l\geq1}$ consisting of selfadjoint contractions in $A$.
These lifts are chosen separately for each element; no multiplicative
lifting of $C$ is needed.

At coordinate $l$, use Proposition~\ref{u:prop:local-moments} with the
finite sets
\[
  H_l=\{c_{1,l},\ldots,c_{l,l}\},\qquad
  E_l=\{x_{1,l},\ldots,x_{l,l}\},
\]
with tolerance $1/l$, and with the same fixed number $R$ at every
coordinate.  Obtain $d_l=d_l^*\in A$, $\norm{d_l}\leq R$, such that
for $k\leq l$,
\begin{align*}
  \norm{[d_l,x_{k,l}]}_{2,T}&<1/l,\\
  \sup_{\rho\in T}|\rho(c_{k,l}d_l)|&\leq3/R^3+1/l,\\
  \sup_{\rho\in T}|\rho(c_{k,l}(d_l^2-1))|&\leq3/R^2+1/l,\\
  \sup_{\rho\in T}\rho(d_l^4)&\leq3+1/l.
\end{align*}
The local proposition permits arbitrary finite subsets of $A$, so their
dependence on $l$ causes no change in its hypotheses.  The common bound
$R$ makes $(d_l)_l$ a bounded sequence.  Let $d$ be its image in $D$.

For a fixed $k$, the displayed estimates hold for all $l\geq k$.
Taking ultralimits, using Equation~\eqref{u:eq:limit-trace-norm} for the
commutators and the definition of limit traces for the moment tests,
gives
\[
  \norm{[d,x_k]}_{2,\Lambda}=0,
  \qquad |\tau(c_kd)|\leq3/R^3,
  \qquad |\tau(c_k(d^2-1))|\leq3/R^2,
  \qquad \tau(d^4)\leq3
\]
for every $\tau\in\Lambda$.  Since the uniform tracial $2$-norm
detects zero in $D$, $[d,x_k]=0$.  Norm continuity and linearity give
$d\in C'$.  For completeness, the relevant continuity estimates are
\[
\begin{split}
  \norm{[d,x-y]}_{2,\Lambda}&\leq2R\norm{x-y},\\
  |\tau((b-c)d)|&\leq R\norm{b-c},\\
  |\tau((b-c)(d^2-1))|&\leq(R^2+1)\norm{b-c}.
\end{split}
\]
They extend the two moment bounds from $(c_k)_k$ to the entire unit
ball of $C$, uniformly over $\Lambda$.  Nuclearity was used in obtaining
the local assertion in $A$; this transfer places no nuclearity
hypothesis on $C$.
\end{proof}

\subsection{Commuting rows and their characteristic functions}

\begin{proposition}[A conditional central limit construction]
\label{u:prop:conditional-clt}
For every integer $N\geq1$ there exist mutually commuting selfadjoint
elements $d_{N,1},\ldots,d_{N,N}\in D\cap A'$, each of norm at most
$N$, with the following properties.  Put
\[
  C_{N,j}=C^*(A,d_{N,1},\ldots,d_{N,j-1}),
  \qquad
  s_N=\frac{1}{\sqrt N}\sum_{j=1}^N d_{N,j}.
\]
For every contraction $b\in C_{N,j}$ and $\tau\in\Lambda$,
\begin{equation}
\label{u:eq:row-moments}
  |\tau(bd_{N,j})|\leq\frac{3}{N^3},\qquad
  |\tau(b(d_{N,j}^2-1))|\leq\frac{3}{N^2},\qquad
  \tau(d_{N,j}^4)\leq3.
\end{equation}
For every fixed $t\in\R$,
\begin{equation}
\label{u:eq:uniform-characteristic-functions}
  \lim_{N\to\infty}
  \sup_{\substack{\tau\in\Lambda\\a\in A,\ \norm{a}\leq1}}
  \left|\tau(ae^{its_N})-e^{-t^2/2}\tau(a)\right|=0.
\end{equation}
\end{proposition}

\begin{proof}
Fix $N$.  At the $j$th step, apply
Proposition~\ref{u:prop:relative-moments} with $R=N$ to $C_{N,j}$.
This algebra is unital and norm separable, since $A$ is separable and
only finitely many elements have been adjoined.  The resulting
$d_{N,j}$ commutes exactly with $A$ and its predecessors and satisfies
Equation~\eqref{u:eq:row-moments}.  This constructs each row separately.

We prove the asserted convergence quantitatively.  Fix $t\in\R$,
$\tau\in\Lambda$, and a contraction $a\in A$.  Define
\[
  P_0=1,\qquad
  P_j=\prod_{k=1}^j e^{itd_{N,k}/\sqrt N},\qquad
  u_j=\tau(aP_j),\qquad
  q_N=1-\frac{t^2}{2N}.
\]
At step $j$, the coefficient $b=aP_{j-1}$ belongs to $C_{N,j}$ and
has norm at most one.  Write $d=d_{N,j}$.  Scalar Taylor expansion
with its integral remainder gives
\[
  |e^{ix}-1-ix+x^2/2|\leq |x|^3/6
  \qquad(x\in\R).
\]
Functional calculus, followed by the tracial inequality
$|\tau(by)|\leq\tau(|y|)$ for a contraction $b$, therefore yields
\begin{align*}
 &\left|\tau(be^{itd/\sqrt N})-\tau(b)
       -\frac{it}{\sqrt N}\tau(bd)
       +\frac{t^2}{2N}\tau(bd^2)\right|\\
 &\hspace{35mm}\leq
       \frac{|t|^3}{6N^{3/2}}\tau(|d|^3)
       \leq\frac{3^{3/4}|t|^3}{6N^{3/2}}.
\end{align*}
The last inequality is H\"older's inequality for the probability
measure associated with $d$ and the tracial state $\tau$, since
$\tau(|d|^3)\leq\tau(d^4)^{3/4}$.  Inserting the two conditional
moment bounds from Equation~\eqref{u:eq:row-moments} proves, for every
$1\leq j\leq N$,
\begin{equation}
\label{u:eq:clt-recurrence}
\begin{split}
  |u_j-q_Nu_{j-1}|&\leq\epsilon_N(t),\\
  \epsilon_N(t)&=
    \frac{3^{3/4}|t|^3}{6N^{3/2}}
    +\frac{3|t|}{N^{7/2}}
    +\frac{3t^2}{2N^3}.
\end{split}
\end{equation}
Every bound here holds for all $\tau\in\Lambda$ and all contractions
$a\in A$, with the same right-hand side.

For fixed $t$, one has $|q_N|\leq1$ once $N\geq t^2/4$.
Telescoping the recurrence gives
\[
  u_N-q_N^Nu_0
    =\sum_{j=1}^N q_N^{N-j}(u_j-q_Nu_{j-1}),
  \qquad
  |u_N-q_N^N\tau(a)|\leq N\epsilon_N(t).
\]
Since the elements in a row commute, $P_N=e^{its_N}$.
Consequently
\begin{equation}
\label{u:eq:quantitative-clt}
\begin{split}
 &\sup_{\substack{\tau\in\Lambda\\a\in A,\ \norm{a}\leq1}}
   |\tau(ae^{its_N})-e^{-t^2/2}\tau(a)|\\
 &\qquad\leq
   \frac{3^{3/4}|t|^3}{6\sqrt N}
   +\frac{3|t|}{N^{5/2}}+\frac{3t^2}{2N^2}
   +\left|\left(1-\frac{t^2}{2N}\right)^N-e^{-t^2/2}\right|
   \longrightarrow0.
\end{split}
\end{equation}
This proves Equation~\eqref{u:eq:uniform-characteristic-functions}.
The coefficients $aP_{j-1}$ are exactly why moment bounds against the
whole earlier algebra were required.  No independence of the variables
or equality of their distributions has been assumed.
\end{proof}

Each $s_N$ is a single element of $D\cap A'$.  The available estimate
$\norm{s_N}\leq N^{3/2}$ is sufficient here: the family $(s_N)_N$ is
not used as a bounded sequence representing an element of an
ultrapower.

\subsection{Uniform weak convergence and Gaussian partitions}

Let $\gamma$ be the standard Gaussian probability measure on $\R$,
so that
\[
  d\gamma(x)=(2\pi)^{-1/2}e^{-x^2/2}\,dx,
  \qquad \int_{\R}e^{itx}\,d\gamma(x)=e^{-t^2/2}.
\]

\begin{lemma}[Uniform weak convergence with a fixed coefficient]
\label{u:lem:uniform-weak-convergence}
For every fixed positive contraction $a\in A$ and every fixed bounded
continuous function $f:\R\to\C$,
\begin{equation}
\label{u:eq:uniform-weak-convergence}
  \lim_{N\to\infty}\sup_{\tau\in\Lambda}
  \left|\tau(af(s_N))-\tau(a)\int_{\R}f\,d\gamma\right|=0.
\end{equation}
\end{lemma}

\begin{proof}
Suppose the conclusion fails.  Then $\norm{f}_\infty>0$, and there
are $\delta>0$, integers $N_k\to\infty$, and traces
$\tau_k\in\Lambda$ such that
\begin{equation}
\label{u:eq:weak-failure}
  \left|\tau_k(af(s_{N_k}))
       -\tau_k(a)\int_{\R}f\,d\gamma\right|\geq\delta
  \qquad(k\geq1).
\end{equation}
The positive functional $x\mapsto\tau_k(a^{1/2}xa^{1/2})$ has norm
$\tau_k(a)$ and agrees with $x\mapsto\tau_k(ax)$ by traciality.
Thus the left side of Equation~\eqref{u:eq:weak-failure} is at most
$2\norm{f}_\infty\tau_k(a)$, and in particular
\[
  \tau_k(a)\geq c:=\frac{\delta}{2\norm{f}_\infty}>0.
\]
On the compact spectrum of $s_{N_k}$, the formula
\[
  h\longmapsto
  \frac{\tau_k(ah(s_{N_k}))}{\tau_k(a)}
\]
defines a positive unital functional.  Its representing probability
measure, viewed as a measure $\mu_k$ on $\R$, satisfies
\[
  \int_{\R}h\,d\mu_k
     =\frac{\tau_k(ah(s_{N_k}))}{\tau_k(a)}
\]
for every bounded continuous $h$ on $\R$.  By
Equation~\eqref{u:eq:uniform-characteristic-functions}, for every fixed
$t\in\R$,
\[
  \left|\int_{\R}e^{itx}\,d\mu_k(x)-e^{-t^2/2}\right|
  \leq\frac{1}{c}
  \sup_{\tau\in\Lambda}
  |\tau(ae^{its_{N_k}})-e^{-t^2/2}\tau(a)|
  \longrightarrow0.
\]
The classical L\'evy continuity theorem for characteristic functions
\cite[Theorem~1.1]{Dobler} now implies $\mu_k\to\gamma$ weakly:
the limiting characteristic
function is that of a probability measure and is continuous at zero.
Hence $\int f\,d\mu_k\to\int f\,d\gamma$.  Multiplying by
$0<\tau_k(a)\leq1$ contradicts Equation~\eqref{u:eq:weak-failure}.
This argument permits the testing trace to vary with $k$, which proves
the asserted uniformity.  It also accommodates the varying spectra of
$s_{N_k}$; no common spectral bound is needed for the continuity
theorem.
\end{proof}

\begin{theorem}[Uniform property Gamma]
\label{u:thm:uniform-gamma}
Let $A$ be a simple, separable, unital, infinite-dimensional nuclear
$C^*$-algebra with nonempty tracial-state space and strict comparison
as in Definition~\ref{u:def:strict-comparison}.  For every integer
$m\geq1$ there exist pairwise orthogonal projections
$p_1,\ldots,p_m\in D\cap A'$ such that
\begin{equation}
\label{u:eq:gamma-projections}
  \sum_{i=1}^m p_i=1,
  \qquad
  \tau(ap_i)=\frac{\tau(a)}{m}
  \quad(a\in A,\ \tau\in\Lambda,\ 1\leq i\leq m).
\end{equation}
In particular, $A$ has uniform property $\Gamma$.
\end{theorem}

\begin{proof}
For $m=1$, take $p_1=1$.  Suppose $m\geq2$, and choose the Gaussian
quantiles
\[
  -\infty=q_0<q_1<\cdots<q_{m-1}<q_m=+\infty
\]
so that the consecutive intervals $I_i$ between $q_{i-1}$ and $q_i$
each have $\gamma$-measure $1/m$.  The assignment of the finite
endpoints is immaterial because $\gamma$ has no atoms.

Fix a finite set $H$ of positive contractions in $A$ and a tolerance
$\varepsilon>0$.  Choose $\theta>0$ so small that
\[
  2\theta<\varepsilon,
  \qquad \frac{17\theta}{16}<\varepsilon^2.
\]
Choose pairwise disjoint open intervals about the finitely many
quantiles $q_1,\ldots,q_{m-1}$ whose union $U$ has
$\gamma(U)<\theta$.  Interpolating linearly across these intervals
gives continuous functions $f_1,\ldots,f_m:\R\to[0,1]$ with
\[
  \sum_{i=1}^m f_i=1,
  \qquad
  f_i=1_{I_i}\ \text{on }\R\setminus U.
\]
At a transition around $q_i$, only $f_i$ and $f_{i+1}$ vary, with one
decreasing from one to zero and the other its complement.  In
particular,
\begin{equation}
\label{u:eq:gaussian-partition-errors}
  \left|\int f_i\,d\gamma-\frac1m\right|<\theta,
  \qquad
  \int(f_i-f_i^2)^2\,d\gamma<\frac{\theta}{16}.
\end{equation}
The second estimate uses $(u-u^2)^2\leq1/16$ for $0\leq u\leq1$,
and the integrand vanishes outside $U$.

Keep these functions fixed.  Apply
Lemma~\ref{u:lem:uniform-weak-convergence} to every pair $(a,f_i)$
with $a\in H$, and to $(1,(f_i-f_i^2)^2)$ for each $i$.
For all sufficiently large $N$, all of these finitely many
convergence errors are less than $\theta$.  Set $v_i=f_i(s_N)$.
These are positive contractions in $D\cap A'$ and satisfy
$\sum_i v_i=1$.  Equation~\eqref{u:eq:gaussian-partition-errors} gives
\begin{align*}
  \sup_{\tau\in\Lambda}
       |\tau(av_i)-\tau(a)/m|
     &<2\theta<\varepsilon \qquad(a\in H),\\
  \norm{v_i-v_i^2}_{2,\Lambda}^2
     &=\sup_{\tau\in\Lambda}
          \tau\big((f_i-f_i^2)^2(s_N)\big)
       <\frac{17\theta}{16}<\varepsilon^2.
\end{align*}
Thus arbitrarily accurate finite sets of trace-splitting and
projection tests can be met by $m$ positive contractions that already
commute with $A$ and sum to one.

Choose a countable norm-dense family $(a_k)_{k\geq1}$ in the positive
unit ball of $A$.  Apply Lemma~\ref{u:lem:tracial-diagonal} to the
$m$ variables constrained to be positive contractions, with the
countable tests
\begin{align*}
  \norm{p_i^2-p_i}_{2,\Lambda}&=0
       &&(1\leq i\leq m),\\
  \norm*{\sum_{i=1}^m p_i-1}_{2,\Lambda}&=0,\\
  \norm{[p_i,a_k]}_{2,\Lambda}&=0
       &&(1\leq i\leq m,\ k\geq1),\\
  \sup_{\tau\in\Lambda}
       |\tau(a_kp_i-a_k/m)|&=0
       &&(1\leq i\leq m,\ k\geq1).
\end{align*}
The preceding finite construction verifies the hypotheses of that
lemma.  In this application the diagonal variables are only the
$m$ positive contractions $v_i$; the elements $s_N$ are not diagonal
variables.  In particular all required lifts lie in fixed norm balls.
The sum-to-one equation is included among the tests, regardless of
which positive contractive lifts of a trial tuple are chosen.

The resulting positive contractions $p_i$ are projections, commute
with $A$, and sum to one, because the uniform tracial $2$-norm
detects zero in $D$ and commutation extends by norm density and
linearity.  They are pairwise orthogonal: for each $i$,
\[
  0=p_i(1-p_i)p_i
    =\sum_{j\ne i}p_ip_jp_i
    =\sum_{j\ne i}(p_jp_i)^*(p_jp_i),
\]
and positivity forces $p_jp_i=0$ for $j\ne i$.  The trace equations
extend from the dense family to every positive contraction, uniformly
over $\Lambda$, and then by scaling and linearity to every element
of $A$.  This proves Equation~\eqref{u:eq:gamma-projections}, which is
uniform property $\Gamma$ as in Definition~\ref{u:def:uniform-gamma}.
\end{proof}

\subsection{Jiang--Su absorption and the traceless case}

\begin{proof}[Proof of Theorem~\ref{u:thm:main}]
If $T(A)\ne\varnothing$, Theorem~\ref{u:thm:uniform-gamma} gives
uniform property $\Gamma$.  Lemma~\ref{u:lem:comparison-convention}
converts the assumed stabilized comparison to the convention used in
Theorem~\ref{u:thm:gamma-criterion}.  That theorem therefore gives
$A\cong A\otimes_{\min}\Zalg$.

Suppose now that $T(A)=\varnothing$.  The comparison convention in
the hypothesis says that every positive element of the stabilization
is Cuntz subequivalent to every nonzero positive element.  We show
directly that every nonzero hereditary subalgebra of $A$ contains an
infinite projection.

Given $0\ne b\in A_+$, the stipulated comparison gives
$1\precsim b$.  The comparison witnesses can be taken in $A$:
view $1$ and $b$ in the corner determined by
$e=1_A\otimes e_{11}$, and compress every stabilized witness $z_n$
to $x_n=ez_ne$.  Then
\[
  x_n^*b x_n=e z_n^*b z_n e\longrightarrow e,
\]
which is $x_n^*b x_n\to1$ after identifying this corner with $A$.
Choose $x\in A$ for which $t=x^*b x$ is invertible and put
\[
  w=b^{1/2}x t^{-1/2}.
\]
Then $w^*w=1$, while
\[
  p=ww^*=b^{1/2}x t^{-1}x^*b^{1/2}\in\Her(b)
\]
is a projection Murray--von Neumann equivalent to $1$.

There are orthogonal nonzero positive elements $b,c\in A$: take
nonnegative continuous functions with disjoint supports near two
distinct spectral points of a nonscalar selfadjoint element.
For this $b$, the projection $p\in\Her(b)$ above annihilates $c$,
so $p\ne1$.  Thus $1$ is an infinite projection.  For any nonzero
hereditary subalgebra $E\subset A$, choose $0\ne h\in E_+$ and
repeat the construction to obtain a projection $p_h\in\Her(h)
\subset E$ equivalent to $1$.  It is infinite inside $E$: an
isometry $w_h$ with $w_h^*w_h=1$ and $w_hw_h^*=p_h$ identifies
$A$ with $p_hAp_h\subset E$, and carries a proper subprojection
equivalent to $1$ to a proper subprojection equivalent to $p_h$.
We have proved the hereditary-subalgebra definition of pure
infiniteness for the simple algebra $A$.

Kirchberg--Phillips' absorption theorem
\cite[Theorem~3.14]{KP} now gives $A\otimes_{\min}\Oinf\cong A$.
The algebra $\Oinf$ is strongly self-absorbing
\cite[Examples~1.14(ii)]{TW}, and Winter's theorem gives
$\Oinf\otimes_{\min}\Zalg\cong\Oinf$
\cite[Theorem~3.1]{WinterSSA}.  Associativity of the minimal tensor
product therefore yields
\[
\begin{split}
  A\otimes_{\min}\Zalg
    &\cong(A\otimes_{\min}\Oinf)\otimes_{\min}\Zalg\\
    &\cong A\otimes_{\min}(\Oinf\otimes_{\min}\Zalg)\\
    &\cong A\otimes_{\min}\Oinf
     \cong A.
\end{split}
\]
The absorption theorems used here require no Universal Coefficient
Theorem hypothesis.  This completes both trace-space cases.
\end{proof}
